%% ma: L12-B12-0019
%% lop: 12
%% chuong: 4
%% bai: 12
%% dang: 12.12.2
%% loai: TLN
%% muc: VD
%% dap_an: "5,49"
%% nguon: "(NBV)-ĐỀ SỐ 15-MỖI NGÀY 1 ĐỀ THI 2025.pdf (D15, MỖI NGÀY 1 ĐỀ - Đề số 15), Phần III Câu 4"
%% kien_thuc: "Bài 12 (tích phân của hàm số có chứa dấu giá trị tuyệt đối, quãng đường từ vận tốc), Bài 11 (nguyên hàm); xác định parabol từ đỉnh và điểm thuộc đồ thị"
%% trang_thai: chinh-thuc
%% ghi_chu: "Cùng dạng 12.12.2 (quãng đường từ hàm vận tốc). Đáp án gốc 5,49 đúng (sympy: 56 căn 7 chia 27 = 5,4875). Đề gốc đồ thị cắt Oy tại 3 (đọc từ hình), bản kho đã ghi rõ trong đề. Hình dựng lại: parabol v = 3t^2 - 8t + 3."
\begin{ex}
Một vật chuyển động thẳng. Khi vật bắt đầu qua điểm $O$, người ta bắt đầu quan sát, đo đạc chuyển động của vật thì thấy vận tốc tức thời của vật tại thời điểm $t$ (s) là $v(t)$ (m/s), hàm vận tốc $v(t)$ có đồ thị phụ thuộc thời gian là một phần của đường parabol đỉnh $I\left(\dfrac43;-\dfrac73\right)$ như hình bên (với $t\ge0$ và mốc tính thời gian là lúc bắt đầu quan sát chuyển động; đồ thị cắt trục tung tại điểm có tung độ $3$). Tính tổng quãng đường vật đi được từ lúc đầu đến lúc vật trở lại điểm $O$ lần cuối cùng (kết quả tính bằng mét và làm tròn đến hàng phần trăm).
\begin{center}
\begin{tikzpicture}[scale=.9]
  \draw[truc] (-.3,0)--(3.6,0) node[below]{$t$};
  \draw[truc] (0,-1.7)--(0,2.4) node[left]{$v$};
  \draw[thick,domain=0:3,samples=80,smooth] plot (\x,{(3*\x*\x-8*\x+3)*.5});
  \draw[phu] (1.3333,0)--(1.3333,-1.1667) (0,-1.1667)--(1.3333,-1.1667);
  \fill (1.3333,-1.1667) circle (1.3pt) node[below right,font=\small] {$I$};
  \node[below left,font=\small] at (0,0) {$O$};
  \node[left,font=\small] at (0,1.5) {$3$};
  \node[above,font=\small] at (1.3333,.05) {$\frac43$};
  \node[left,font=\small] at (0,-1.1667) {$-\frac73$};\end{tikzpicture}
\end{center}
\shortans{5,49}
\loigiai{Parabol $v(t)=at^2+bt+c$ có $c=3$, đỉnh $I\left(\dfrac43;-\dfrac73\right)$: $-\dfrac b{2a}=\dfrac43$ và $a\cdot\dfrac{16}9+b\cdot\dfrac43+3=-\dfrac73$, suy ra $a=3$, $b=-8$, $v(t)=3t^2-8t+3$.

Vị trí $s(t)=\displaystyle\int v\,\mathrm dt=t^3-4t^2+3t$ ($s(0)=0$). $s(t)=0\iff t=0$ hoặc $t=1$ hoặc $t=3$, nên lần cuối vật về $O$ là lúc $t=3$ s.

Quãng đường: $S=\displaystyle\int_0^3|3t^2-8t+3|\,\mathrm dt=\dfrac{56\sqrt7}{27}\approx5{,}49$ m.}
\end{ex}
